\section{Kết luận và kế hoạch Giai đoạn 2}

\subsection{Tổng kết}

Giai đoạn~1 đã thiết lập nền móng cho một nền tảng phần mềm quản lý và khai thác
dữ liệu ảnh y khoa phục vụ chẩn đoán xuất huyết não trong các bối cảnh bệnh viện
khác nhau. Kiến thức lâm sàng nền tảng đã được khảo sát, qua đó làm rõ các yêu
cầu mà quy trình đọc ảnh đặt ra cho hệ thống lưu trữ, đặc biệt là việc phải bảo
toàn dữ liệu điểm ảnh gốc để bác sĩ có thể tuỳ ý chuyển đổi cửa sổ hiển thị khi
đọc.

Chuẩn DICOM đã được trình bày như một ràng buộc kiến trúc chứ không phải một chi
tiết cài đặt: nó quy định cách dữ liệu được tiếp nhận, những thuộc tính phải
được bảo toàn để hiển thị chính xác, và những thuộc tính phải được loại bỏ trước
khi chia sẻ. Các giải pháp phần mềm hiện có đã được khảo sát và phân tích, cho
thấy một khoảng trống rõ ràng giữa các hệ thống thương mại hoàn chỉnh nhưng đắt
đỏ và khép kín, với các thành phần mã nguồn mở tuy miễn phí nhưng đòi hỏi công
sức tích hợp đáng kể.

Khung pháp lý Việt Nam, chủ yếu là Luật Bảo vệ dữ liệu cá nhân năm 2025 (kế thừa
và thay thế Nghị định 13/2023/NĐ-CP), Luật Khám bệnh, chữa bệnh và Thông tư
13/2025/TT-BYT về hồ sơ bệnh án điện tử, xác lập dữ liệu sức khoẻ là dữ liệu cá nhân nhạy cảm và điều
chỉnh việc chuyển dữ liệu ra nước ngoài chặt chẽ tới mức mô hình lưu trữ cục bộ
trước, đồng bộ đám mây sau trở thành lựa chọn mặc định hợp lý chứ không đơn
thuần là một ưu tiên thận trọng.

Một mô hình chung đã được xây dựng trước khi đi vào từng bệnh viện: luồng xử lý
chín bước của một ca chụp, nguyên tắc ghi cục bộ trước, các cơ chế độ bền D1--D7
bảo đảm ca chụp đã xác nhận không bị mất, và bốn tầng lưu trữ từ kho nóng tại chỗ
tới kho lạnh ngoài cơ sở. Ba mẫu bối cảnh bệnh viện đã được đặc tả, và ba kiến trúc
tương ứng được đề xuất như các cấu hình của mô hình chung đó, khác nhau ở nơi giữ
bản gốc, cách hiện thực các tầng lưu trữ và hành vi khi có sự cố; mô hình liên viện
trục--nan mở rộng kiến trúc C cho một mạng lưới cơ sở tuyến huyện. Kiến trúc B đã
được kiểm chứng bằng bản mẫu cho phần tiếp nhận, hiển thị và kích hoạt theo sự kiện.

\subsection{Giai đoạn 2}

\begin{enumerate}
    \item \textbf{Hoàn thiện nền tảng.} Hiện thực bộ điều phối kèm khử định
          danh, cơ chế thử lại và kiểm soát tải; hoàn thiện hàng đợi đồng bộ bền
vững mô tả tại Mục~\ref{sec:flow}.
\item \textbf{Lưu trữ bền và phân tầng.} Hiện thực các tầng T1--T3 trên kho
      tương thích S3, chính sách vòng đời và nạp trước ca cũ, cùng các cơ chế
      D3--D7 (kiểm tra toàn vẹn định kỳ, sao lưu 3-2-1 mã hoá, khoá ghi một lần,
      diễn tập khôi phục) theo Mục~\ref{sec:durability} và Mục~\ref{sec:tiering}.
\item \textbf{Tích hợp HIS/RIS.} Nhận chỉ định qua HL7 v2 hoặc FHIR, cung cấp
      Modality Worklist khi có thể, và chức năng đối chiếu định danh theo tinh
      thần IHE PIR (Mục~\ref{sec:his}).
    \item \textbf{An toàn thông tin.} Bổ sung xác thực và phân quyền theo vai
          trò, mã hoá kênh truyền và dữ liệu lưu trữ, nhật ký truy vết chỉ ghi
          thêm, theo đúng các yêu cầu tại Mục~\ref{sec:legal}.
    \item \textbf{Hoàn thiện giao diện đọc ảnh.} Tuỳ biến trình xem cho phù hợp
          với quy trình đọc ảnh sọ não, bao gồm các cửa sổ hiển thị đặt sẵn mô
          tả tại Mục~\ref{sec:windowing}.
    \item \textbf{Đánh giá.} Đo thời gian truyền thực tế và lợi ích của nén không mất
      mát, tải lũy tiến HTJ2K trên đường truyền chậm để thay các ước lượng ở
      Hình~\ref{fig:transfer}; đo RPO và RTO thực tế qua diễn tập sự cố; đo độ trễ đầu-cuối từ lúc thiết bị gửi ảnh tới lúc
          bác sĩ xem được, đo thông lượng, và đặc tả hành vi hệ thống dưới các
          kịch bản lỗi mà nguyên tắc~P1 hướng tới kiểm soát, đặc biệt là kịch
          bản mất kết nối kéo dài.
    \item \textbf{Triển khai thử nghiệm.} Với điều kiện có thoả thuận hợp tác
          phù hợp, triển khai thử tại một cơ sở y tế nhằm kiểm chứng giả định về
          hạ tầng và quy trình đã nêu ở Mục~\ref{sec:contexts}.
\end{enumerate}

Rủi ro chính đối với Giai đoạn~2 là khả năng tiếp cận một cơ sở y tế thực tế để
triển khai thử nghiệm, yếu tố mà việc kiểm chứng các giả định kiến trúc phụ
thuộc vào. Vì vậy, việc kết nối với các cơ sở y tế cần được bắt đầu sớm, và kế
hoạch đánh giá cần được xây dựng sao cho có thể chuyển về phương án kiểm thử
trên dữ liệu mô phỏng nếu không thu xếp được triển khai thực tế.
